\subsection{Exact low-degree coefficient comparisons in every weight}

Fix $d\geq0$.  Let $\mathcal A^{[d]}\subset\mathcal A_{\mathbb Z}$
be the full subcategory in which the coefficient comparison is
an isomorphism in degree zero and has vanishing positive
cohomology through degree $d$, in \emph{every} weight and simplex.
Let $\mathcal C^{[d]}=\mathcal C\cap\mathcal A^{[d]}$, where
$\mathcal C$ is the full subcategory of A-006.

\begin{theorem}[A-S09-24: Proved; both main questions reduce to global low-degree exactness]
\label{AS0924}
The inclusions
\[
 \mathcal A^{[d]}\hookrightarrow\mathcal A_{\mathbb Z},
 \qquad
 \mathcal C^{[d]}\hookrightarrow\mathcal A_{\mathbb Z}
 \tag{S09.115}
\]
induce equivalences on the groupoids obtained by inverting
every morphism.  In particular they preserve and reflect
terminality of those localizations.
The common-source assertion for $\mathcal A_{\mathbb Z}$ is
equivalent to the assertion with sources and targets all internal
to either displayed full subcategory.

The subcategories are proper.  More precisely,
$\mathcal A^{[d]}$ is proper in $\mathcal A_{\mathbb Z}$,
$\mathcal C^{[d]}$ is proper in $\mathcal C$, and
$\mathcal C^{[d]}$ is proper in $\mathcal A^{[d]}$.
The equivalences and properness statements also hold with the
support condition $q\leq w$, with termwise flatness, and with
both conditions together.
\end{theorem}

\begin{proof}
For $N\geq1$ define the weight-tail replacement by
\[
 (T_NA)^{w,*}=
 \begin{cases}
 S^w[0],&w<N,\\
 A^{w,*},&w\geq N.
 \end{cases}
 \tag{S09.116}
\]
Here all operations below the cutoff are the coefficient
operations; a mixed product whose weight is at least $N$
uses the structural coefficient map to $A$ and its product.
The analogous rule defines every divided power.  One precise
finite-limit description, which verifies every relation at once,
is
\[
        T_NA=A\mathop{\times}_{A_{<N}}S_{<N},
       \qquad A_{<N}=A/\!\bigoplus_{w\geq N}A^{w,*}.
 \tag{S09.117}
\]
The high-weight ideal is a differential PD ideal, since weights
are nonnegative and divided powers multiply the weight.
Thus (S09.117) is a pullback of actual augmented strict PD
cochain algebras in every simplex.  Its componentwise
identifications with (S09.116) give the stated operations and
all simplicial maps.

There is a natural strict map
$t_{N,A}:T_NA\to A$: it is the structural coefficient map in
weights below $N$ and the identity above.  It need not be
injective when the original low-weight coefficient map is
noninjective.  For a strict map $f:A\to B$, the functor $T_Nf$
is the identity on the replaced coefficient pieces and is $f$
on the remaining pieces.  The pullback or the displayed formulas
prove functoriality and naturality exactly.

For $M\geq N$ there are actual natural maps
$T_MA\to T_NA$, given by coefficients in weights $N\leq w<M$
and by the identity elsewhere.  In the componentwise models,
\[
                   T_MT_N=T_{\max\{M,N\}},
 \tag{S09.118}
\]
including all products, differentials, divided powers and
coefficient maps.  The transition maps compose to the indicated
transition for the maximum; their composites with $t_N$ are
$t_M$.

The replacement remains admissible, with bounds
$\max\{N,\lambda_A(e)\}$ and
$\max\{N,\mu_A(e)\}$ for every $e$.
Indeed its rows and complexes in the stated high range are
literally those of $A$.  It preserves support and flatness,
since each replaced piece is a free coefficient group in form
degree zero.  If $N\geq\max\{1,\lambda_A(d)\}$, every weight
of (S09.116) has the required comparison through degree $d$.
Thus $T_NA\in\mathcal A^{[d]}$.  Also $T_N$ preserves
$\mathcal A^{[d]}$ for every $N$.

For the second subcategory use the actual functor of A-006 and
put
\[
 U_N=R T_N,\qquad
 \theta_{N,A}=t_{N,A}\,\tau_{T_NA}:U_NA\longrightarrow A.
 \tag{S09.119}
\]
This is a functor on the full admissible category with an actual
natural transformation to its identity, for each fixed $N$.
Choose
\[
             N\geq
        \max\{1,\lambda_A(d+2),\mu_A(d+2)\}.
 \tag{S09.120}
\]
The input $T_NA$ then has the vertical coefficient comparison
through $d+2$ in every weight, and horizontal exactness through
form degree $d+2$ in every weight at least $N$.
For $0<w<N$ its only zero-forms are the coefficient group,
with both face maps the identity.  Therefore
$E_w(T_NA)=\varnothing$: an element with $d_1a=0$ is zero and
cannot have $d_0a=T^{[w]}$ in the free coefficient group.
Every generator of $P_{T_NA}$ has weight at least $N$.
Its coskeleton and the pullback defining $U_NA$ consequently
equal the coefficient algebra in weights below $N$.
Above $N$, apply the actual comparison and bounds of A-006
with the just-established controls for $T_NA$.
They prove the coefficient comparison through degree $d$.
Hence
\[
                         U_NA\in\mathcal C^{[d]}
                   \quad\hbox{for every $N$ in (S09.120)}.
 \tag{S09.121}
\]
No functorial choice of a numerical cutoff has been asserted.
For any finite collection of objects, a single maximum of the
finitely many displayed bounds gives (S09.121) for all of them.

We prove full faithfulness on the localizations explicitly.
The same argument applies first with
$(F_N,\eta_N,\mathcal D)=(T_N,t_N,\mathcal A^{[d]})$ and then
with $(U_N,\theta_N,\mathcal C^{[d]})$.
These are the concrete functors and transformations constructed
above.  A morphism in an all-arrow localization is a finite
word in actual arrows and their formal inverses.  For objects
$X,Y\in\mathcal D$, take any such word $w:X\to Y$ in the
full admissible category.  Choose one $N$ making $F_N$ of
each object occurring in that word an object of $\mathcal D$.
The word
\[
                      [\eta_{N,Y}]\,[F_N(w)]\,
                                      [\eta_{N,X}]^{-1}
 \tag{S09.122}
\]
is entirely a word in morphisms of $\mathcal D$.
Naturality on each arrow, and then on its formal inverse,
identifies its image with $w$.  This proves fullness.

For faithfulness, let two words $u,v$ internal to $\mathcal D$
become equal in the full localization.  In the presentation
of that groupoid, equality is generated by insertion or
removal of identity arrows, composition relations for two
actual composable arrows, and cancellation of an arrow with
its formal inverse.  An equality is a finite sequence of
these generating changes.  It therefore involves finitely
many objects and arrows.  Choose one $N$ for that entire
finite witness.  Functoriality sends every displayed relation
to the same valid relation inside $\mathcal D$, so
$F_N(u)=F_N(v)$ in its localization.  The naturality identity
(S09.122) for $u$ and $v$ now holds inside that localization:
their original objects and all their $F_N$ images lie in
$\mathcal D$, and it is a full subcategory.  Conjugating by
the two endpoint components proves $u=v$ there.  This proves
faithfulness, with the equality witness and all comparison
maps specified.

Finally every admissible $A$ receives the actual arrow
$\eta_{N,A}:F_NA\to A$ from an object of $\mathcal D$ for a
sufficiently large $N$.  It becomes an isomorphism after
localization.  This proves essential surjectivity and hence
both equivalences in (S09.115).  They preserve and reflect
equivalence to the one-object groupoid with only its identity
arrow, which proves the terminality assertion.

The common-source assertions also use actual maps.  If all
admissible pairs have common sources and $X,Y\in\mathcal D$,
take $D\to X,Y$ and choose $N$ with $F_ND\in\mathcal D$.
The composites
$F_ND\xrightarrow{\eta_{N,D}}D\to X,Y$
are a common source internal to the full subcategory.
Conversely, for arbitrary $A,B$ choose one $N$ making
$F_NA,F_NB$ internal.  A common source there, composed with
$\eta_{N,A},\eta_{N,B}$, is a common admissible source for
$A,B$.  Every support and flatness restriction is preserved
by all these actual replacements and compositions.

We finish with explicit properness examples.  Let $C_0=R(G)$,
an object of $\mathcal C$.  Its weight-zero algebra is exactly
$\mathbb Z$ in form degree zero: the standard $G$, the free
disk first skeleton, its matching coskeleton and the pullback
all have this property.  Weight-zero projection is therefore
an actual simplicial augmented PD-DGA map
\[
                              \alpha:C_0\longrightarrow\mathbb Z.
 \tag{S09.123}
\]
This is the augmentation to the ground ring; no coefficient
retraction to $S$ is used.
Let $M=K(\mathbb Z,2)$, placed in a fixed weight $b\geq1$
and form degree zero, with zero vertical differential.
Form the strict square-zero extension $F=C_0\ltimes M$,
using the action through $\alpha$:
\[
 (c,m)(c',m')=(cc',\alpha(c)m'+\alpha(c')m),\qquad
 d(c,m)=(dc,0),\qquad M^2=0.
 \tag{S09.124}
\]
On the augmentation ideal set $\gamma_1(m)=m$ and
$\gamma_n(m)=0$ for $n\geq2$, and use the PD addition rule
with the old operations on $C_0$.  The mixed terms for
$n\geq2$ vanish because the positive-weight ideal acts
trivially on $M$.  These formulas verify the PD and
differential identities directly and give split strict maps
$C_0\to F\to C_0$.

The module $M$ vanishes in simplex degrees zero and one,
so $F$ has the same first skeleton and belongs to
$\mathcal C$.  It differs from $C_0$ only in weight $b$;
increasing each eventual cutoff beyond $b$ proves
admissibility.  Its terms are flat and satisfy $q\leq w$
when those of $C_0$ do.  In simplex two and weight $b$,
its $H^0$ contains the new direct summand $M_2=\mathbb Z$,
whereas the coefficient map lands in the old summand.
Thus $F\notin\mathcal A^{[d]}$ for every $d\geq0$.
This proves properness of $\mathcal A^{[d]}$ in the full
category and of $\mathcal C^{[d]}$ within $\mathcal C$.

For properness within $\mathcal A^{[d]}$, use the split
example at the end of A-005: adjoin to $G$ a free even-to-odd
identity disk on $D(1)$ in forms zero and one and weight one.
Its integral disk contraction gives the vertical coefficient
comparison in every simplex and weight, so it lies in every
$\mathcal A^{[d]}$.  Its simplex-zero weight-one one-form
is nonzero, which excludes a first-skeleton isomorphism with
any $P(E)$.  It is outside $\mathcal C$, and hence outside
$\mathcal C^{[d]}$.  The same example is flat and supported,
completing all the stated properness assertions.
\end{proof}
